% =========================
% Chapter 9
% =========================
\chapter{Future Work}
\label{ch:future_work}

The limitations of Chapter~\ref{ch:limitations} map directly onto a staged plan.
The items below are ordered roughly by leverage: each earlier item either unblocks
or de-risks the ones after it.

\section{Deepening the Human-in-the-Loop Navigation Design}
\label{sec:fw_hints}

The exploration design casts the user as a \emph{direction oracle}, and the core
of that interaction is now live: at a multi-door fork the system announces the
doors, asks, follows a spoken ``left''/``right''/``ahead''/``behind'', and falls
back to the door nearest straight ahead when the user does not know
(Section~\ref{sec:multi_door}). The next milestones deepen it:

\begin{itemize}[nosep]
  \item \textbf{Hints in the initial command.} Accept a direction embedded in the
        task utterance itself (``the kitchen is to the left''), so the very first
        scan can already be biased before any fork is reached. The parser's
        direction vocabulary and the bearing mathematics are shared with the fork
        dialogue; what is missing is carrying the hint into the task context.
  \item \textbf{Richer door references.} Beyond four coarse directions: ordinal
        and comparative references (``the second door on the right'', ``the far
        one''), and a spoken correction channel after guidance has started
        (``not that one'').
  \item \textbf{Come-from-door exclusion.} After a doorway transit, the entry
        door sits roughly $180^\circ$ from the entry heading, which the compass
        provides nearly for free. Deprioritizing (not banning) it eliminates the
        two-room ping-pong that is currently the main failure mode of longer
        journeys --- and removes the most common reason the fork question has to
        be asked at all.
\end{itemize}

\section{Exploration Memory and Bounded Autonomy}
\label{sec:fw_memory}

Beyond hints lies coarse memory --- deliberately heuristic, never a map:

\begin{itemize}[nosep]
  \item \textbf{Tried-door memory and backtracking.} Bucket each door as
        unvisited, entry, or exhausted; prefer unvisited, backtrack through the
        entry door at dead ends, and avoid exhausted branches --- depth-first
        search over an unknown room graph, with a breadcrumb stack of entry
        headings enabling multi-level backtracking. The compass prerequisite is
        already wired.
  \item \textbf{Room cap with a graceful hand-off.} Bound how far exploration
        can wander: after a configured number of rooms, do not quit silently ---
        let the user decide: ``I'm having trouble finding the kitchen. Do you
        want to stop, or keep going?'' The user's ``stop'' is already honored
        globally; the cap makes the check-in proactive.
\end{itemize}

\section{Perception Robustness}
\label{sec:fw_perception}

\begin{itemize}[nosep]
  \item \textbf{Hard-negative fine-tuning.} The highest-value perception fix:
        30--50 photographs of the deployment environment's walls, curtains, and
        wardrobes, added to the door datasets with empty labels (``nothing
        here''), and a short fine-tune of both door models. This teaches the
        models that these surfaces are \emph{not} doors instead of fencing
        around the weakness with gates --- and would allow some gates to be
        relaxed, recovering recall.
  \item \textbf{Distance calibration.} A single ground-truth measurement per
        device class locks the one-shot calibration constant; per-device camera
        parameters (true field of view) would remove the nominal-FOV assumption
        entirely.
  \item \textbf{Vocabulary growth.} Additional destination types (hallway,
        stairs, exit) require indicator classes COCO lacks; a small fine-tune
        for architectural classes --- stairs, elevators, exit signs --- would
        extend both navigation targets and obstacle awareness. Automated tests
        for the perception layer's pure parts (edge-density and IoU logic
        against synthetic frames) would close the current live-only validation
        gap.
\end{itemize}

\section{User Studies}
\label{sec:fw_studies}

The decisive evaluation is a formal study with blind and low-vision participants,
run under a safety protocol (sighted spotter, cane in use, staged environments
first). Beyond task-level metrics --- completion rates, times, interventions ---
the study should target the claims the design is actually built on: that
task-scoped activation lowers cognitive load relative to continuous narration,
that the speech discipline earns trust, and that step-count and clock-free
directional phrasing are actionable for the population. Instrument-supported
measures (NASA-TLX for workload, structured interviews for trust) would
complement the task metrics. The findings should be expected to reshape the
guidance phrasing at least as much as the original live trials reshaped
perception.

\section{Platform Evolution}
\label{sec:fw_platform}

\begin{itemize}[nosep]
  \item \textbf{Reducing network dependence.} Progressively move latency-critical
        pieces client-side: on-device detection for the obstacle watchdog (the
        one component where reaction time is safety-relevant), offline TTS, or a
        fully on-device mode on capable phones --- keeping the server path for
        the heavy navigation stack.
  \item \textbf{Hardened deployment.} A persistent HTTPS host in place of the
        developer tunnel, retention policies for captured audio, debug tracing
        behind a runtime flag, and load validation of the per-session isolation
        the architecture already provides.
  \item \textbf{Multilingual voice.} Whisper is inherently multilingual and gTTS
        supports many languages; the work is in the parser's template and
        synonym vocabularies --- with Arabic the natural first target for the
        project's context.
\end{itemize}